\DocumentMetadata{
  pdfversion=2.0,
  pdfstandard=ua-2,
  lang=en-US,
  tagging=on,
  tagging-setup={math/setup=mathml-SE,tabsorder=structure}
}
\documentclass[11pt,letterpaper]{article}
\usepackage[margin=0.68in,includefoot]{geometry}
\usepackage{newcomputermodern}
\usepackage{amsmath,amscd,mathtools}
\usepackage{tikz,adjustbox}
\providecommand{\square}{\mdlgwhtsquare}
\usepackage[hidelinks]{hyperref}
\hypersetup{
  pdftitle={Logic PhD Exam, May 2023},
  pdfauthor={University of Florida Department of Mathematics},
  pdfsubject={Graduate mathematics examination},
  pdfdisplaydoctitle=true,
  bookmarksopen=true
}
\setcounter{secnumdepth}{0}
\setlength{\parindent}{0pt}
\setlength{\parskip}{0.28em}
\setlength{\emergencystretch}{3em}
\setlength{\leftmargini}{2.15em}
\setlength{\leftmarginii}{2.65em}
\setlength{\leftmarginiii}{3em}
\pagestyle{empty}
\raggedbottom
\allowdisplaybreaks
\NewTaggingSocketPlug{block/list/label}{examproblem}{%
  \tagstructbegin{tag=Lbl}\tagstructend%
  \tagstructbegin{tag=\LBody}%
  \tagstructbegin{tag=\ProblemHeadingTag,title-o={\ProblemHeadingText}}%
  \tagmcbegin{tag=\ProblemHeadingTag,actualtext={\ProblemHeadingText}}%
  #2%
  \tagmcend%
  \tagstructend%
}
\newenvironment{examproblems}[2]{%
  \def\ProblemHeadingTag{#2}%
  \AssignTaggingSocketPlug{block/list/label}{examproblem}%
  \begin{enumerate}%
  \setcounter{enumi}{#1}%
  \renewcommand{\labelenumi}{\textbf{\theenumi.}}%
  \setlength{\topsep}{0.35em}%
  \setlength{\partopsep}{0pt}%
  \setlength{\itemsep}{0.48em}%
  \setlength{\parsep}{0.18em}%
}{\end{enumerate}}
\newenvironment{examsubparts}{%
  \AssignTaggingSocketPlug{block/list/label}{default}%
  \begin{enumerate}%
  \setlength{\topsep}{0.18em}%
  \setlength{\partopsep}{0pt}%
  \setlength{\itemsep}{0.16em}%
  \setlength{\parsep}{0.1em}%
}{\end{enumerate}}
\newenvironment{examinstructions}{%
  \par\begingroup\itshape
}{%
  \par\endgroup\vspace{0.18em}
}
\makeatletter
\renewcommand\subsection{\@startsection{subsection}{2}{0pt}%
  {-1.25ex plus -.25ex minus -.1ex}%
  {0.55ex plus .1ex}%
  {\normalfont\large\bfseries}}
\renewcommand{\@maketitle}{%
  \newpage\null\vskip -1em%
  {\footnotesize\bfseries
    \href{https://www.ufl.edu/}{University of Florida}, \href{https://math.ufl.edu/}{Department of Mathematics}\par}%
  \vskip -0.5em%
  \tagstructbegin{tag=H1,title={Logic PhD Exam, May 2023}}%
  \tagpdfparaOff%
  \tagmcbegin{tag=H1}%
    {\LARGE\bfseries\href{https://gma.math.ufl.edu/exams/logic-phd/}{Logic PhD Exam}, May 2023\par}%
  \tagmcend%
  \tagpdfparaOn%
  \tagstructend%
  \vskip 0.25em%
  \tagmcbegin{artifact}%
    \hrule height 1pt%
  \tagmcend%
  \vskip 1em%
}
\makeatother
\title{Logic PhD Exam, May 2023}
\author{}
\date{}
\begin{document}
\maketitle
\thispagestyle{empty}
\begin{examinstructions}
Solve 5 problems of the following; at least one from each section.
\end{examinstructions}
\subsection{A. Set Theory.}
\begin{examproblems}{0}{H3}
\def\ProblemHeadingText{Problem 1.}
\item
State Martin’s Axiom for \(\aleph_1\) and prove that it implies the negation of the Continuum Hypothesis.
\def\ProblemHeadingText{Problem 2.}
\item
Give the definition of a standard Borel space. Sketch the proof that any two uncountable standard Borel spaces are Borel isomorphic.
\def\ProblemHeadingText{Problem 3.}
\item
Let~\(\kappa\) be an uncountable cardinal. Show that there is an ultrafilter on~\(\kappa\) containing no sets of cardinality smaller than~\(\kappa\).
\end{examproblems}
\subsection{B. Computability.}
\begin{examproblems}{0}{H3}
\def\ProblemHeadingText{Problem 1.}
\item
State two different common definitions of the class of computable functions and sketch the proof that they give the same class.
\def\ProblemHeadingText{Problem 2.}
\item
Define the many-one reducibility and sketch the proof that there is a many-one degree strictly between~\(0\) and \(0'\).
\def\ProblemHeadingText{Problem 3.}
\item
Prove that every infinite computably enumerable subset of the natural numbers contains an infinite computable subset.
\end{examproblems}
\subsection{C. Model theory.}
\begin{examproblems}{0}{H3}
\def\ProblemHeadingText{Problem 1.}
\item
What is quantifier elimination for theories? Give an example of a theory which has it, and prove that it has it. Give an example of theory which does not have it, and prove that it does not have it.
\def\ProblemHeadingText{Problem 2.}
\item
State Łoś’s theorem and prove it.
\def\ProblemHeadingText{Problem 3.}
\item
Give an example of a relational Fraïssé class which has the disjoint amalgamation property, and an example of a relational Fraïssé class which does not have it.
\end{examproblems}
\end{document}
